\newif\ifglobal

%%%%%%%%%%%%%%%%%%%%%%%
%%% Compile to view %%%
%%%%%%%%%%%%%%%%%%%%%%%

%%%%%%%%%%%%%%%%%%%%%%%%%%%%%%%%%%%%%%%%%%%%%%%%%%%
%%% Readme:                                     %%%
%%% https://www.overleaf.com/read/bwdjvfjksvjy/ %%%
%%%%%%%%%%%%%%%%%%%%%%%%%%%%%%%%%%%%%%%%%%%%%%%%%%%

% >>> M?: marks a module setting number or important notice

% >>> M4: Set {./relative-path} to external files for this module <<<
\def \modulefiles {./45-I2S}
% To add an external file, use:
% (Replace '---' with actual filename)
% '\input{\modulefiles/tables/---.tex}' for tables
% '\includegraphics{\modulefiles/figures/---}' for figures
% '\subfile{\modulefiles/---}' for register files

% >>> M1: This module is in global project? <<<
% 'YES' - uncomment; 'NO' - comment out
\globaltrue

\ifglobal
    \documentclass[main\_\_EN.tex]{subfiles}
   
\else
    \documentclass[openany, 10pt]{book}

    % >>> M2: In ./00-shared/preamble-trm-module, also find
    %         settings for visibility of todo notes, watermark <<<
    \usepackage{preamble-trm-module}

    % >>> M3: Temporary labels in English,
    %         you can add your own labels there <<<
    \myexternaldocument{./00-shared/temp-labels-en}

\fi %\ifglobal


\setcounter{secnumdepth}{4}


% >>> M5: If this module needs any updates in
% './00-shared/preamble-trm-module.sty'
% (include additional package, etc.), contact your module PM
% or create the file './00-shared/changelog.tex' make the
% necessary updates yourself and document them in this file <<<

% Make text left-aligned and allow latex to break pages naturally, comment out for full justification
\raggedright
\raggedbottom


\begin{document}


% Sets language into which pre-defined bits of text are translated
\selectlanguage{english}
\setenglishformatting

% Do not delete this block
\ifglobal
% Add the following front matter if
% this module is outside of global project
\else
    \InputIfFileExists{readme}{}{}
    \clearpage
    \listoftodos
    \listofdonetodos
    \tableofcontents
    %\listoftables
    %\listoffigures
\fi


%%%%%%%%%%%%%%%%%%%%%%%%%
%%% TRM Module Begins %%%
%%%%%%%%%%%%%%%%%%%%%%%%%

\hypertarget{i2s}{}
\chapter{I2S Controller (I2S)}
\label{mod:i2s}

\section{Overview}
\label{sec:i2s-overview}

\chipname{} has two I2S interfaces, I2S\regindex{0} and I2S\regindex{1}, which provide flexible data communication interfaces for multimedia applications, especially digital audio applications.

The I2S standard bus defines three signals: bit clock (BCK), word select (WS), and serial data (SD). A basic I2S data bus has one master and one slave, whose roles remain unchanged throughout communication. The \chipname{} I2S module has independent TX and RX units for high-performance communication.

\begin{tiplisting}
The information in this chapter applies to both I2S\regindex{0} and I2S\regindex{1}. Unless otherwise stated, I2S and I2S\regindex{n} in this chapter refer to both units.

The following features are available only on I2S\regindex{0}:
\begin{itemize}
    \item PDM TX and PDM RX modes, including raw PDM data transmission and reception, the PCM-to-PDM converter, and the PDM-to-PCM converter
    \item Dedicated PDM data lines I2S\regindex{0}O\_SD1\_out, I2S\regindex{0}I\_SD1\_in, I2S\regindex{0}I\_SD2\_in, and I2S\regindex{0}I\_SD3\_in
    \item RX data recombination by channel into multiple DMA interfaces
\end{itemize}
The corresponding register fields remain readable and writable in I2S\regindex{1}, but have no effect and must not be configured by software. Except for these features, I2S\regindex{0} and I2S\regindex{1} are identical.
\end{tiplisting}

\section{Terminology}
\label{sec:i2s-term}

To better illustrate the functionality of I2S, the following terms are used in this chapter. 

\begin{longtable}[c]{ R{5cm} p{10cm} }
\textbf{Master mode}     & As a master, I2S\regindex{n} drives BCK/WS and transmits data to or receives data from a slave. \\ 
\textbf{Slave mode}      & As a slave, I2S\regindex{n} is driven by BCK/WS and receives data from or transmits data to a master. \\
\textbf{Full-duplex} & Transmit and receive use separate data lines and can run at the same time.\\
\textbf{Half-duplex} & Transmit and receive share a data line and cannot run at the same time.\\
\textbf{A-law and $\mu$-law} & A-law and $\mu$-law are compression/decompression algorithms in digital pulse code modulated (PCM) non-uniform quantization, which can effectively improve the signal-to-quantization noise ratio. \\
\textbf{TDM RX mode} & In this mode, pulse code modulated (PCM) data is received utilizing time division multiplexing (TDM). The signal lines include BCK, WS, and SD. Data from 16 channels at most can be received. TDM Philips standard, TDM MSB alignment standard, and TDM PCM standard are supported in this mode, depending on the configuration.\\
\textbf{PDM RX mode (I2S\regindex{0} only)} & In this mode, pulse density modulation (PDM) data is received. The signal lines include WS and SD (DATA). The data format is configurable and supports the PDM standard.\\
\textbf{TDM TX mode} & In this mode, pulse code modulated (PCM) data is sent in a way of time division multiplexing (TDM). The signal lines include BCK, WS, and SD (DATA). Data up to 16 channels can be sent. TDM Philips standard, TDM MSB alignment standard, and TDM PCM standard are supported in this mode, depending on the configuration.\\
\textbf{PDM TX mode (I2S\regindex{0} only)} & In this mode, pulse density modulation (PDM) data is transmitted. The signal lines include WS and SD (DATA). The data format is configurable and supports the PDM standard.\\
\textbf{PCM-to-PDM data format converter (I2S\regindex{0} only)} & A data format converter that converts PCM data to PDM data (hereinafter the PCM-to-PDM converter). It can be enabled in PDM TX master mode. When enabled, I2S\regindex{0} fetches PCM data from memory through GDMA, converts it to PDM data, and transmits it.\\
\textbf{PDM-to-PCM data format converter (I2S\regindex{0} only)} & A data format converter that converts PDM data to PCM data (hereinafter the PDM-to-PCM converter). It can be enabled in PDM RX master or slave mode. When enabled, I2S\regindex{0} receives PDM data, converts it to PCM data, and stores it in memory through GDMA.\\
\end{longtable}

\section{Features}
\label{sec:i2s-features}

The I2S module has the following features:

\begin{itemize}
    \item Master mode and slave mode
    \item Full-duplex and half-duplex communications
    \item Separate TX and RX units that can work independently or simultaneously
    \item A variety of audio standards supported: 
        \begin{itemize}
            \item TDM Philips standard
            \item TDM MSB alignment standard
            \item TDM PCM standard
            \item PDM standard (I2S\regindex{0} only)
        \end{itemize}
    \item Various TX/RX modes supported:
        \begin{itemize}
        \item TDM TX mode, up to 16 channels supported
        \item TDM RX mode, up to 16 channels supported
        \item PDM TX mode (I2S\regindex{0} only)
        \begin{itemize}
            \item Raw PDM data transmission
            \item PCM-to-PDM data format conversion, up to two channels supported
        \end{itemize}
        \item PDM RX mode (I2S\regindex{0} only)
        \begin{itemize}
            \item Raw PDM data reception
            \item PDM-to-PCM conversion, supporting up to eight channels
        \end{itemize}
        \end{itemize}
    \item Configurable clock source with frequencies up to 125 MHz
    \item Configurable high-precision sample clock with a variety of sampling frequencies supported (refer to the note below for more details)
    \item 8/16/24/32-bit data width
    \item Synchronous counter in TX mode
    \item RX data recombination by channel into up to four DMA interfaces (I2S\regindex{0} only)
    \item ETM feature
    \item GDMA
    \item I2S interface interrupts
    \item A-law and $\mu$-law compression/decompression
\end{itemize}

\begin{tiplisting}
In slave mode, the maximum sampling frequency supported by \chipname{} I2S\regindex{n} is limited by the clock-source frequency, data width, and number of channels. For example, with two 32-bit channels, sampling frequencies below 122 kHz are supported, including 8 kHz, 16 kHz, 32 kHz, 44.1 kHz, 48 kHz, 88.2 kHz, and 96 kHz. See Section \ref{sec:i2s-tx-rx-clk} \textit{\nameref{sec:i2s-tx-rx-clk}}.
\end{tiplisting}

\section{System Architecture}
\label{sec:i2s-arch-overview}

Figure \ref{fig:i2s-sys-diagram} shows the structure of the \chipname{} I2S module, consisting of:

\begin{itemize}
    \item Independent TX unit
    \item Independent RX unit
    \item Input/output timing adjustment unit (I/O Sync)
    \item Clock divider (Clock Generator)
    \item 64 × 32-bit TX FIFO
    \item 64 × 32-bit RX FIFO
    \item Compress/Decompress units
\end{itemize}

\begin{figure}[H]
    \centering
    \includegraphics[width=0.9\textwidth]{\modulefiles/figures/i2s_arch_20260626__EN}
    \\ Note: The PDM-to-PCM and PCM-to-PDM converters are available only on I2S\regindex{0}.
    \caption{\chipname{} I2S System Diagram}
    \label{fig:i2s-sys-diagram}
\end{figure}

The \chipname{} I2S module supports GDMA access to internal memory. For more information, see Chapter \ref{mod:dma} \textit{\nameref{mod:dma}}. 

Both the TX unit and the RX unit have a 3-line interface that uses a bit clock line (BCK), a word select line (WS), and a serial data line (SD). The SD line of the TX unit is used for data output and the SD line of the RX unit for data input. The BCK and WS signal lines of the TX and RX units are used for data output in master mode, and the BCK and WS signal lines of the TX and RX units are used for data input in slave mode.

The right side of Figure \ref{fig:i2s-sys-diagram} shows the I2S\regindex{n} signal bus. RX and TX signal names follow the pattern I2S\regindex{n}\textcolor{red}{A}\_\textcolor{red}{B}\_\textcolor{red}{C}, for example, I2S\regindex{n}\textcolor{red}{I}\_\textcolor{red}{BCK}\_\textcolor{red}{in}.

\begin{itemize}
    \item ``A" indicates that the signal belongs to the I2S TX or RX unit, which includes:
    \begin{itemize}
        \item ``I": Signal input to/output from the RX unit
        \item ``O": Signal input to/output from the TX unit
    \end{itemize}
    \item ``B" represents the signal function, which includes:
    \begin{itemize}
        \item BCK
        \item WS
        \item SD
    \end{itemize}
    \item ``C" represents the signal direction, which includes:
    \begin{itemize}
        \item ``in": Input signal into the I2S module
        \item ``out": Output signal from the I2S module
    \end{itemize}
\end{itemize}

Table \ref{tab:i2s-signal-descr} provides a detailed description of I2S signals.

\begin{ThreePartTable}
\begin{TableNotes}
    \item[*] All I2S signals must be mapped to chip pins through the GPIO matrix. For more information, see Chapter \ref{mod:iomuxgpio} \textit{\nameref{mod:iomuxgpio}}.
\end{TableNotes}
\begin{longtable}[c]{|p{3cm}|p{1.5cm}|p{8cm}|}
\caption{I2S Signal Description}
\label{tab:i2s-signal-descr}\\\hline
    \rowcolor{lightgray}
    \textbf{Signal} & \textbf{Direction} & \textbf{Function} \\
\hline
\endfirsthead
    \rowcolor{lightgray}
    \textbf{Signal} & \textbf{Direction} & \textbf{Function} \\
\hline
\endhead
\endfoot
\insertTableNotes
\endlastfoot
    I2S\regindex{n}I\_BCK\_in  & Input & In I2S slave mode, inputs the BCK signal for the RX unit. \\
    \hline
    I2S\regindex{n}I\_BCK\_out& Output & In I2S master mode, outputs the BCK signal for the RX unit. \\
    \hline
    I2S\regindex{n}I\_WS\_in & Input & In I2S slave mode, inputs the WS signal for the RX unit. \\
    \hline
    I2S\regindex{n}I\_WS\_out& Output & In I2S master mode, outputs the WS signal for the RX unit. \\\hline
    I2S\regindex{n}I\_SD\_in& Input & Serial input data line for the RX unit.\\\hline
    I2S\regindex{n}O\_SD\_out& Output & Serial output data line for the TX unit.\\\hline
    I2S\regindex{n}O\_BCK\_in& Input & In I2S slave mode, inputs the BCK signal for the TX unit. \\\hline
    I2S\regindex{n}O\_BCK\_out& Output & In I2S master mode, outputs the BCK signal for the TX unit. \\\hline
    I2S\regindex{n}O\_WS\_in& Input & In I2S slave mode, inputs the WS signal for the TX unit. \\\hline
    I2S\regindex{n}O\_WS\_out& Output & In I2S master mode, outputs the WS signal for the TX unit. \\\hline
    I2S\regindex{n}\_MCLK\_in& Input & External clock source in I2S slave mode.\\\hline
    I2S\regindex{n}\_MCLK\_out& Output & Clock source for an external device in I2S master mode.\\\hline
    I2S\regindex{0}I\_SD1\_in& Input & Serial input data line for the RX unit when the PDM-to-PCM converter is enabled.\\\hline
    I2S\regindex{0}I\_SD2\_in& Input & Serial input data line for the RX unit when the PDM-to-PCM converter is enabled.\\\hline
    I2S\regindex{0}I\_SD3\_in& Input & Serial input data line for the RX unit when the PDM-to-PCM converter is enabled.\\\hline
    I2S\regindex{0}O\_SD1\_out& Output & Serial output data line for the TX unit when the PCM-to-PDM converter is enabled in 2-line DAC output format. \\\hline
\end{longtable}
\end{ThreePartTable}

\section{Functional Description}
\label{sec:i2s-func-descr}

\subsection{Audio Standards Supported by I2S} \label{sec:i2s-supported-audio-standards}

\chipname{} I2S\regindex{n} supports the TDM Philips, TDM MSB alignment, TDM PCM, and PDM standards. The PDM standard is available only on I2S\regindex{0}.

Select the needed standard by configuring the following bits:

\begin{itemize}
    \item \hyperref[fielddesc:I2SRXTDMEN]{I2S\_TX/RX\_TDM\_EN}
    \begin{itemize}
        \item 0: Disable TDM mode
        \item 1: Enable TDM mode
    \end{itemize}
    \item \hyperref[fielddesc:I2SRXPDMEN]{I2S\_TX/RX\_PDM\_EN}
    \begin{itemize}
        \item 0: Disable PDM mode
        \item 1: Enable PDM mode
    \end{itemize}    
    \item \hyperref[fielddesc:I2SRXMSBSHIFT]{I2S\_TX/RX\_MSB\_SHIFT}
    \begin{itemize}
        \item 0: WS and SD signals change simultaneously, i.e., enable MSB alignment standard
        \item 1: WS signal changes one BCK clock cycle earlier than SD signal, i.e., enable Philips standard or select PCM standard
    \end{itemize}
\end{itemize}

\begin{tiplisting}
\hyperref[fielddesc:I2SRXTDMEN]{I2S\_TX/RX\_TDM\_EN} and \hyperref[fielddesc:I2SRXPDMEN]{I2S\_TX/RX\_PDM\_EN} must not both be 1 or both be 0. Otherwise, \chipname{} I2S\regindex{n} operates incorrectly in a mode that is neither TDM nor PDM.
\end{tiplisting}

\subsubsection{TDM Philips Standard Mode}

Philips standards require that WS signal changes one BCK clock cycle earlier than SD signal on BCK falling edge, which means that WS signal is valid from one clock cycle before transmitting the first bit of channel data and changes one clock before the end of channel data transfer. SD signal line transmits the most significant bit of audio data first. 

Compared with the basic Philips standard, TDM Philips standard supports multiple channels. See Figure \ref{fig:i2s-tdm-philips-std-timing-diagram}.

\begin{figure}[H]
    \centering
    \includegraphics[width=1.0\textwidth]{\modulefiles/figures/i2s_tdm_phillips_mode.png}
    \caption{TDM Philips Standard Timing Diagram}
    \label{fig:i2s-tdm-philips-std-timing-diagram}
\end{figure}

\subsubsection{TDM MSB Alignment Standard Mode}

MSB alignment standards require that WS and SD signals change simultaneously on the falling edge of BCK. The WS signal is valid until the end of the channel data transfer. The SD signal line transmits the most significant bit of audio data first.

Compared with the basic MSB alignment standard, TDM MSB alignment standard supports multiple channels. See Figure \ref{fig:i2s-tdm-msb-align-std-timing-diagram}.

\begin{figure}[H]
    \centering
    \includegraphics[width=1.0\textwidth]{\modulefiles/figures/i2s_tdm_msb_mode.png}
    \caption{TDM MSB Alignment Standard Timing Diagram}
    \label{fig:i2s-tdm-msb-align-std-timing-diagram}
\end{figure}

\subsubsection{TDM PCM Standard Mode}

Short frame synchronization under the PCM standards requires that the WS signal changes one BCK clock cycle earlier than the SD signal on the falling edge of BCK, which means that the WS signal becomes valid one clock cycle before transferring the first bit of channel data and remains unchanged in this BCK clock cycle. SD signal line first transmits the most significant bit of audio data. 

Compared with the basic PCM standard, TDM PCM standard supports multiple channels. See Figure \ref{fig:i2s-tdm-pcm-std-timing-diagram}.

\begin{figure}[H]
    \centering
    \includegraphics[width=1.0\textwidth]{\modulefiles/figures/i2s_tdm_pcm_mode.png}
    \caption{TDM PCM Standard Timing Diagram}
    \label{fig:i2s-tdm-pcm-std-timing-diagram}
\end{figure}

\subsubsection{PDM Standard Mode}

Under PDM standard, WS signal changes continuously during data transmission. The low-level and high-level of this signal indicates the left channel and right channel respectively. WS and SD signals change simultaneously. See Figure \ref{fig:i2s-pdm-std-timing-diagram}. 

\begin{figure}[H]
    \centering
    \includegraphics[width=0.8\textwidth]{\modulefiles/figures/i2s_pdm_mode.pdf}
    \caption{PDM Standard Timing Diagram}
    \label{fig:i2s-pdm-std-timing-diagram}
\end{figure}

\subsection{I2S Master/Slave Mode}

The \chipname{} I2S module can operate as a master or a slave in half-duplex and full-duplex communications, depending on the configuration of  \hyperref[fielddesc:I2SRXSLAVEMOD]{I2S\_RX\_SLAVE\_MOD} and  \hyperref[fielddesc:I2STXSLAVEMOD]{I2S\_TX\_SLAVE\_MOD}.

\begin{itemize}
    \item \hyperref[fielddesc:I2STXSLAVEMOD]{I2S\_TX\_SLAVE\_MOD} 
    \begin{itemize}
        \item 0: Master TX mode
        \item 1: Slave TX mode
    \end{itemize}
    \item \hyperref[fielddesc:I2SRXSLAVEMOD]{I2S\_RX\_SLAVE\_MOD}
    \begin{itemize}
        \item 0: Master RX mode
        \item 1: Slave RX mode
    \end{itemize}
\end{itemize}

\subsubsection{Master/Slave TX Mode} \label{sec:i2s-master-slave-tx-mode}

\begin{itemize}
    \item I2S works as a master transmitter:
    \begin{itemize}
        \item Set \hyperref[fielddesc:I2STXSTART]{I2S\_TX\_START} to start transmitting data. When this bit is set, the TX unit keeps driving the clock signal and serial data.
        \item If \hyperref[fielddesc:I2STXSTOPEN]{I2S\_TX\_STOP\_EN} is set and all the data in FIFO is transmitted, the master stops transmitting data and clock signals.
        \item If \hyperref[fielddesc:I2STXSTOPEN]{I2S\_TX\_STOP\_EN} is cleared and all the data in FIFO is transmitted, meanwhile no new data is filled into FIFO, the TX unit keeps transmitting the last data frame and clock signal.
        \item Master stops transmitting data when \hyperref[fielddesc:I2STXSTART]{I2S\_TX\_START} is cleared.
    \end{itemize}
    \item I2S works as a slave transmitter:
     \begin{itemize}
        \item Set \hyperref[fielddesc:I2STXSTART]{I2S\_TX\_START}. Wait for the master BCK clock to enable a transmit operation.
        \item If \hyperref[fielddesc:I2STXSTOPEN]{I2S\_TX\_STOP\_EN} is set and all the data in FIFO is transmitted, then the slave keeps transmitting zeros, till the master stops providing BCK signal.
        \item If \hyperref[fielddesc:I2STXSTOPEN]{I2S\_TX\_STOP\_EN} is cleared and all the data in FIFO is transmitted, meanwhile no new data is filled into FIFO, the TX unit keeps transmitting the last data frame. 
        \item If \hyperref[fielddesc:I2STXSTART]{I2S\_TX\_START} is cleared, slave keeps transmitting zeros till the master stops providing BCK clock signal.
    \end{itemize}
\end{itemize}

\subsubsection{Master/Slave RX Mode}\label{sec:i2s-master-slave-rx-mode}

\begin{itemize}
    \item I2S works as a master receiver:
    \begin{itemize}
        \item Set \hyperref[fielddesc:I2SRXSTART]{I2S\_RX\_START} to start receiving data. When this bit is set, the RX unit keeps outputting clock signal and sampling input data.
        \item Configure \hyperref[fielddesc:I2SRXSTOPMODE]{I2S\_RX\_STOP\_MODE} to control the suspension of data reception:
        \begin{itemize}
            \item 0: The RX unit only suspends data reception when \hyperref[fielddesc:I2SRXSTART]{I2S\_RX\_START} is cleared.
            \item 1: The RX unit suspends data reception when \hyperref[fielddesc:I2SRXSTART]{I2S\_RX\_START} is cleared or the number of received bytes reaches the receive length configured in \hyperref[regdesc:I2SRXEOFNUMREG]{I2S\_RXEOF\_NUM\_REG}. If \hyperref[fielddesc:I2SRXSTART]{I2S\_RX\_START} remains set, set \hyperref[fielddesc:I2SRXRESET]{I2S\_RX\_RESET} to restart reception.
            \item 2: The RX unit suspends data reception when \hyperref[fielddesc:I2SRXSTART]{I2S\_RX\_START} is cleared or the trigger condition of \hyperref[int:I2SRXDONEINT]{I2S\regindex{n}\_RX\_DONE\_INT} is met. The condition is configured by \hyperref[fielddesc:I2SRXDONEMODE]{I2S\_RX\_DONE\_MODE}; see Section \ref{sec:i2s-interrupts} \textit{\nameref{sec:i2s-interrupts}}. If \hyperref[fielddesc:I2SRXSTART]{I2S\_RX\_START} remains set, set \hyperref[fielddesc:I2SRXRESET]{I2S\_RX\_RESET} to restart reception.
        \end{itemize}
        \item RX unit stops receiving data when \hyperref[fielddesc:I2SRXSTART]{I2S\_RX\_START} is cleared.
    \end{itemize}
    \item I2S works as a slave receiver:
    \begin{itemize}
        \item Set \hyperref[fielddesc:I2SRXSTART]{I2S\_RX\_START}. Wait for the master BCK signal to start receiving data.
        \item Configure \hyperref[fielddesc:I2SRXSTOPMODE]{I2S\_RX\_STOP\_MODE} to control data reception suspension:
        \begin{itemize}
            \item 0: The RX unit only suspends data reception when \hyperref[fielddesc:I2SRXSTART]{I2S\_RX\_START} is cleared.
            \item 1: The RX unit suspends data reception when \hyperref[fielddesc:I2SRXSTART]{I2S\_RX\_START} is cleared or the number of received bytes reaches the receive length configured in \hyperref[regdesc:I2SRXEOFNUMREG]{I2S\_RXEOF\_NUM\_REG}. If \hyperref[fielddesc:I2SRXSTART]{I2S\_RX\_START} remains set, set \hyperref[fielddesc:I2SRXRESET]{I2S\_RX\_RESET} to restart reception.
            \item 2: The RX unit suspends data reception when \hyperref[fielddesc:I2SRXSTART]{I2S\_RX\_START} is cleared or the trigger condition of \hyperref[int:I2SRXDONEINT]{I2S\regindex{n}\_RX\_DONE\_INT} is met. The condition is configured by \hyperref[fielddesc:I2SRXDONEMODE]{I2S\_RX\_DONE\_MODE}; see Section \ref{sec:i2s-interrupts} \textit{\nameref{sec:i2s-interrupts}}. If \hyperref[fielddesc:I2SRXSTART]{I2S\_RX\_START} remains set, set \hyperref[fielddesc:I2SRXRESET]{I2S\_RX\_RESET} to restart reception.
        \end{itemize}
        \item The RX unit stops receiving data when \hyperref[fielddesc:I2SRXSTART]{I2S\_RX\_START} is cleared.
    \end{itemize}
\end{itemize}

\subsection{Transmitting Data}\label{sec:i2s-tx-data}

\vspace{1em}
\begin{tiplisting}
Updating the configuration described in this and subsequent sections requires to set \hyperref[fielddesc:I2STXUPDATE]{I2S\_TX\_UPDATE} accordingly to synchronize registers from APB clock domain to TX clock domain. For more detailed configurations, see Section \ref{sec:i2s-configure-i2s-tx-mode}.
\end{tiplisting}

\vspace{-1em}
In TX mode, I2S first reads data through GDMA and transmits these data out via output signals according to the configured data mode and channel mode.
\subsubsection{Data Format Control}\label{sec:i2s-data-control}
Data format is controlled in the following phases:
\begin{itemize}
    \item Phase I: Read data from memory and write it to TX FIFO.
    \item Phase II: Read the TX data from TX FIFO and convert the data according to the output data mode.
    \item Phase III: Clock out the TX data serially.
\end{itemize}

\paragraph{Channel Valid Data Width}

The bit width of the valid data in each channel is determined by \hyperref[fielddesc:I2STXBITSMOD]{I2S\_TX\_BITS\_MOD} and \hyperref[fielddesc:I2STX24FILLEN]{I2S\_TX\_24\_FILL\_EN}. For details, see the table below.

\begin{table}[H]
    \centering
    \caption{Bit Width of Channel Valid Data}
    \label{tab:i2s-bit-width-chan-valid-data}
    \begin{threeparttable}
    \begin{tabular}{|p{4.2cm}|R{3.1cm}|R{3.5cm}|}
    \hline
    \rowcolor{lightgray}
    \textbf{Channel Valid Data Width} & \textbf{\hyperref[fielddesc:I2STXBITSMOD]{I2S\_TX\_BITS\_MOD}}  & \textbf{\hyperref[fielddesc:I2STX24FILLEN]{I2S\_TX\_24\_FILL\_EN}} \\ \hline
    \multirow{2}*{32} & 31 & x\tnote{*} \\\cline{2-3}
                      & 23 & 1 \\\hline

                24    & 23 & 0 \\\hline

                16     & 15 & x \\\hline
                8     & 7 & x \\\hline
    \end{tabular}
        \begin{tablenotes}
            \item[*] ``x" represents that this value is ignored (applies to all tables in this chapter).
        \end{tablenotes}
    \end{threeparttable}
\end{table}

\paragraph{Channel Valid Data Endianness}

When I2S reads data through GDMA, the data endian under various data width is controlled by \hyperref[fielddesc:I2STXBIGENDIAN]{I2S\_TX\_BIG\_ENDIAN}. Table \ref{tab:i2s-endian-chan-valid-data} shows how \hyperref[fielddesc:I2STXBIGENDIAN]{I2S\_TX\_BIG\_ENDIAN} controls the reading of the data with different valid data widths of the channel.

\begin{table}[H]
    \centering
    \caption{Endian of Channel Valid Data}
    \label{tab:i2s-endian-chan-valid-data}
    \begin{tabular}{|p{4.2cm}|R{2.7cm}|R{4.2cm}|R{3.5cm}|}
    \hline
    \rowcolor{lightgray}
    \textbf{Channel Valid Data Width} & \textbf{Original Data}  & \textbf{Endian of Processed Data}  & \textbf{\hyperref[fielddesc:I2STXBIGENDIAN]{I2S\_TX\_BIG\_ENDIAN}} \\ \hline
    \multirow{2}*{32} & \multirow{2}*{\{B3, B2, B1, B0\}} & \{B3, B2, B1, B0\} & 0 \\\cline{3-4}
                      && \{B0, B1, B2, B3\} & 1 \\\hline

    \multirow{2}*{24} & \multirow{2}*{\{B2, B1, B0\}} & \{B2, B1, B0\} & 0 \\\cline{3-4}
                      && \{B0, B1, B2\} & 1 \\\hline

    \multirow{2}*{16} & \multirow{2}*{\{B1, B0\}} & \{B1, B0\} & 0 \\\cline{3-4}
                      && \{B0, B1\} & 1 \\\hline

                8     & \{B0\} & \{B0\} & x\tnote{} \\\hline
    \end{tabular}
\end{table}

\begin{tiplisting}
B0, B1, B2, B3 each represents an 8-bit data, and the symbol \{\} indicates that the bytes are combined together. For example, \{B3, B2, B1, B0\} represents a 32-bit data, wherein B0 represents bits 0--7, B1 represents bits 8--15, B2 represents bits 16--23, and B3 represents bits 24--31.
\end{tiplisting}

\paragraph{A-law/$\mu$-law Compression and Decompression}

\chipname{} I2S compresses/decompresses the valid data into 32-bit by A-law or by $\mu$-law. If the bit width of valid data is smaller than 32, zeros are filled to the extra high bits of the data to be compressed/decompressed by default.

\begin{tiplisting}
Extra high bits here mean the bits[31: channel valid data width] of the data to be compressed/decompressed.
\end{tiplisting}

Configure \hyperref[fielddesc:I2STXPCMBYPASS]{I2S\_TX\_PCM\_BYPASS}:
\begin{itemize}
\item 0: Compress or decompress the data
\item 1: Do not compress or decompress the data
\end{itemize}
Configure \hyperref[fielddesc:I2STXPCMCONF]{I2S\_TX\_PCM\_CONF}: 
\begin{itemize}
\item 0: Decompress the data using A-law
\item 1: Compress the data using A-law
\item 2: Decompress the data using $\mu$-law
\item 3: Compress the data using $\mu$-law
\end{itemize}

At this point, the first phase of data format control is completed.

\paragraph{Channel TX Data Width}
\label{sec:i2s-tdm-chan-bits}

The TX data width in each channel is determined by \hyperref[fielddesc:I2STXTDMCHANBITS]{I2S\_TX\_TDM\_CHAN\_BITS}.

\begin{itemize}
\item If TX data width in each channel is larger than the valid data width, zeros will be filled to these extra bits. Configure  \hyperref[fielddesc:I2STXLEFTALIGN]{I2S\_TX\_LEFT\_ALIGN}:
\begin{itemize}
\item 0: The valid data is at the lower bits of TX data. Zeros are filled into higher bits of TX data.
\item 1: The valid data is at the higher bits of TX data. Zeros are filled into lower bits of TX data.
\end{itemize}
\item If the TX data width in each channel is smaller than the valid data width, only the lower bits of valid data are sent out, and the higher bits are discarded.
\end{itemize}

At this point, the second phase of data format control is completed.

\paragraph{Channel Data Bit Order}

The data bit order in each channel is controlled by \hyperref[fielddesc:I2STXBITORDER]{I2S\_TX\_BIT\_ORDER}:
\begin{itemize}
\item 0: Not reverse the valid data bit order.
\item 1: Reverse the valid data bit order.
\end{itemize}

At this point, the data format control is completed. The data after format control will be sent sequentially from high to low. Figure \ref{fig:i2s-tx-data-format-control} shows the complete process of TX data format control.

\begin{figure}[H]
    \centering
    \includegraphics[width=0.7\textwidth]{\modulefiles/figures/i2s_tx_data_config.png}
    \caption{TX Data Format Control}
    \label{fig:i2s-tx-data-format-control}
\end{figure}

\subsubsection{Channel Mode Control}

\chipname{} I2S supports both TDM TX mode and PDM TX mode. Set \hyperref[fielddesc:I2STXTDMEN]{I2S\_TX\_TDM\_EN} to enable TDM TX mode, or set \hyperref[fielddesc:I2STXPDMEN]{I2S\_TX\_PDM\_EN} to enable PDM TX mode. PDM TX mode is available only on I2S\regindex{0}.

\begin{tiplisting}
\hyperref[fielddesc:I2STXTDMEN]{I2S\_TX\_TDM\_EN} and \hyperref[fielddesc:I2STXPDMEN]{I2S\_TX\_PDM\_EN} must not be cleared or set simultaneously. 
\end{tiplisting}
\vspace{-1em}

\paragraph{TDM TX Mode}\label{sec:i2s-chan-control-tdm-mode} 

In TDM TX mode, I2S supports up to 16 channels to transmit data. The total number of TX channels in use is controlled by \hyperref[fielddesc:I2STXTDMTOTCHANNUM]{I2S\_TX\_TDM\_TOT\_CHAN\_NUM}. For example, if \hyperref[fielddesc:I2STXTDMTOTCHANNUM]{I2S\_TX\_TDM\_TOT\_CHAN\_NUM} is set to 5, six channels in total (channels 0--5) will be used to transmit data. See Figure \ref{fig:i2s-tdm-chan-control}.

\vspace{-1em}
\begin{tiplisting}
Most stereo I2S codecs can be controlled by setting the I2S module into 2-channel mode under TDM standard.
\end{tiplisting}

In these TX channels, if \hyperref[fielddesc:I2STXTDMCHAN0EN]{I2S\_TX\_TDM\_CHAN{\textcolor{red}{n}}\_EN} is set to:

\begin{itemize}
    \item 1: This channel transmits the channel data out.
    \item 0: The TX data to be sent by this channel is controlled by \hyperref[fielddesc:I2STXCHANEQUAL]{I2S\_TX\_CHAN\_EQUAL}:
    \begin{itemize}
        \item 1: The data of the previous channel is sent out.
        \item 0: The data stored in \hyperref[fielddesc:I2SSINGLEDATA]{I2S\_SINGLE\_DATA} is sent out.
    \end{itemize}
\end{itemize}

In TDM TX master mode, WS signal is controlled by \hyperref[fielddesc:I2STXWSIDLEPOL]{I2S\_TX\_WS\_IDLE\_POL} and \hyperref[fielddesc:I2STXTDMWSWIDTH]{I2S\_TX\_TDM\_WS\_WIDTH}:
\begin{itemize}
    \item \hyperref[fielddesc:I2STXWSIDLEPOL]{I2S\_TX\_WS\_IDLE\_POL}: The default level of WS signal.
    \item \hyperref[fielddesc:I2STXTDMWSWIDTH]{I2S\_TX\_TDM\_WS\_WIDTH}: Configures the duration for which WS remains at its default level during transmission. The duration is (\hyperref[fielddesc:I2STXTDMWSWIDTH]{I2S\_TX\_TDM\_WS\_WIDTH} + 1) × T\_BCK.
\end{itemize}

(\hyperref[fielddesc:I2STXHALFSAMPLEBITS]{I2S\_TX\_HALF\_SAMPLE\_BITS} + 1) × 2 is equal to the BCK cycles in one WS period.

\textbf{TDM Channel Configuration Example}

In this example, the register configuration is as follows:

\begin{itemize}
    \item \hyperref[fielddesc:I2STXTDMTOTCHANNUM]{I2S\_TX\_TDM\_TOT\_CHAN\_NUM} = 5, i.e., channels 0--5 are used to transmit data.
    \item \hyperref[fielddesc:I2STXCHANEQUAL]{I2S\_TX\_CHAN\_EQUAL} = 1, i.e., data of the previous channel will be transmitted if the \hyperref[fielddesc:I2STXTDMCHAN0EN]{I2S\_TX\_TDM\_CHAN\regindex{n}\_EN} (\regindex{n} = 0--5) is cleared.
    \item \hyperref[fielddesc:I2STXTDMCHAN0EN]{I2S\_TX\_TDM\_CHAN0/2/5\_EN} = 1, i.e., these channels transmit their channel data out.
    \item \hyperref[fielddesc:I2STXTDMCHAN0EN]{I2S\_TX\_TDM\_CHAN1/3/4\_EN} = 0, i.e., these channels transmit the previous channel's data out.
\end{itemize}

Once the configuration is done, data is transmitted as follows.

\begin{figure}[H]
    \centering
    \includegraphics[width=1.0\textwidth]{\modulefiles/figures/i2s_tdm_channel_config.pdf}
    \caption{TDM Channel Control}
    \label{fig:i2s-tdm-chan-control}
\end{figure}

\paragraph{PDM TX Mode} 

This subsection applies only to I2S\regindex{0}. In this mode, I2S\regindex{0} supports both raw PDM data transmission and PCM-to-PDM conversion.

In PDM TX mode, fetching data through GDMA is controlled by \hyperref[fielddesc:I2STXMONO]{I2S\_TX\_MONO} and \hyperref[fielddesc:I2STXMONOFSTVLD]{I2S\_TX\_MONO\_FST\_VLD}. See Table \ref{tab:i2s-data-fetch-control-pdm-mode}. Configure the two bits according to the data stored in memory, be it the single-channel or dual-channel data. 
\begin{table}[H]
    \centering
    \caption{Data-Fetching Control in PDM Mode}
    \label{tab:i2s-data-fetch-control-pdm-mode}
    \begin{tabular}{|p{6.6cm}|l|r|r|}
    \hline
    \rowcolor{lightgray}
    \textbf{Data-Fetching Control Option} &\textbf{Mode} &\textbf{\hyperref[fielddesc:I2STXMONO]{I2S\_TX\_MONO}}  & \textbf{\hyperref[fielddesc:I2STXMONOFSTVLD]{I2S\_TX\_MONO\_FST\_VLD}} \\ \hline
                        Post data-fetching request to GDMA at any edge of WS signal & Stereo mode & 0 & x        \\ \hline
                        Post data-fetching request to GDMA only at the second half period of WS signal  & Mono mode & 1 & 0   \\ \hline
                        Post data-fetching request to GDMA only at the first half period of WS signal & Mono mode & 1 & 1      \\ \hline
    \end{tabular}
\end{table}

When I2S is in PDM TX master mode, the default level of WS signal is controlled by \hyperref[fielddesc:I2STXWSIDLEPOL]{I2S\_TX\_WS\_IDLE\_POL}, and the WS signal frequency is half of the BCK signal frequency. The configuration of WS signal is similar to that of BCK signal. Please refer to Section \ref{sec:i2s-tx-rx-clk} \textit{\nameref{sec:i2s-tx-rx-clk}} and Figure \ref{fig:i2s-pdm-chan-control-example} \textit{\nameref{fig:i2s-pdm-chan-control-example}}.

When \textbf{transmitting raw PDM data}, the I2S channel mode is controlled by  \hyperref[fielddesc:I2STXCHANMOD]{I2S\_TX\_CHAN\_MOD} and  \hyperref[fielddesc:I2STXWSIDLEPOL]{I2S\_TX\_WS\_IDLE\_POL}. See the table below. 

\begin{table}[H]
    \centering
    \caption{I2S Channel Control for Raw PDM Data}
    \label{tab:i2s-chan-control-pdm-mode}
    \begin{threeparttable}
    \begin{tabular}{|p{2.4cm}|L{4.7cm}|L{4.7cm}|R{1.3cm}|R{1.4cm}|}
    \hline
    \rowcolor{lightgray}
    \textbf{Channel Control Option} & \textbf{Left Channel} &\textbf{Right Channel} & \textbf{Mode Control Field}\tnote{1} & \textbf{Channel Select Bit}\tnote{2}\\ \hline
    Stereo mode & Transmit the left channel data & Transmit the right channel data &0 & x \\\hline
                         \multirow{8}*{Mono mode}  & Transmit the left channel data & Transmit the left channel data & 1 & 0         \\ \cline{2-5}
                                           & Transmit the right channel data & Transmit the right channel data& 1 & 1\\ \cline{2-5}
                                           & Transmit the right channel data & Transmit the right channel data & 2 & 0 \\ \cline{2-5}
                                           & Transmit the left channel data & Transmit the left channel data &2 & 1 \\ \cline{2-5}
                                           & Transmit the value of ``single"\tnote{3} & Transmit the right channel data & 3 & 0 \\ \cline{2-5}
                                           & Transmit the left channel data & Transmit the value of ``single" &3 & 1 \\ \cline{2-5}
                                           & Transmit the left channel data & Transmit the value of ``single"&4 & 0 \\ \cline{2-5}
                                           & Transmit the value of ``single" & Transmit the right channel data & 4 & 1 \\ \hline
    \end{tabular}
            \begin{tablenotes}
            \item[1] \hyperref[fielddesc:I2STXCHANMOD]{I2S\_TX\_CHAN\_MOD}
            \item[2] \hyperref[fielddesc:I2STXWSIDLEPOL]{I2S\_TX\_WS\_IDLE\_POL}
            \item[3] The ``single" value is equal to the value of  \hyperref[fielddesc:I2SSINGLEDATA]{I2S\_SINGLE\_DATA}.
        \end{tablenotes}

    \end{threeparttable}
\end{table}

If the \textbf{PCM-to-PDM converter is enabled}, the PCM data through GDMA is converted to PDM data and then output in PDM signal format. Configure \hyperref[fielddesc:I2SPCM2PDMCONVEN]{I2S\_PCM2PDM\_CONV\_EN} to enable the converter. The register configuration for the PCM-to-PDM converter is as follows:

\begin{itemize}
    \item Configure 1-line PDM output format or 1-/2-line DAC output mode as the table below:

\begin{table}[H]
    \centering
    \caption{I2S PCM-to-PDM Data Output}
    \label{tab:i2s-pcm-pdm-output-mode}
    \begin{tabular}{|p{3.9cm}|R{5.2cm}|R{5cm}|}
    \hline
    \rowcolor{lightgray}
    \textbf{Channel Output Format}&\textbf{\hyperref[fielddesc:I2STXPDMDACMODEEN]{I2S\_TX\_PDM\_DAC\_MODE\_EN}} & \textbf{\hyperref[fielddesc:I2STXPDMDAC2OUTEN]{I2S\_TX\_PDM\_DAC\_2OUT\_EN}}\\ \hline
    1-line PDM output format\hyperref[note-1]{$^1$}& 0 & x \\ \hline
    1-line DAC output format\hyperref[note-2]{$^2$} &1   & 0  \\ \hline 
    2-line DAC output format & 1 & 1 \\ \hline
    \end{tabular}
\end{table}

\begin{tiplisting}
\vspace{-2em}
\begin{enumerate}
    \item \label{note-1} In PDM output format, SD data of two channels is sent out in one WS period.
    \item \label{note-2} In DAC output format, SD data of one channel is sent out in one WS period.
    \item 1-line PDM/DAC output format uses I2S\regindex{0}O\_SD\_out as the data output line.
    \item 2-line DAC output format uses I2S\regindex{0}O\_SD\_out and I2S\regindex{0}O\_SD1\_out as data output lines.
    \item In terms of application, the PDM output format is targeted at applications that require clock signals and can decode PDM format codecs, while the DAC output format is targeted at applications that do not rely on clock signals, do not have PDM format codecs, and can recover analog waveforms directly through low-pass filtering.
\end{enumerate}
\end{tiplisting}
    
    \item Configure sampling frequency and upsampling rate as below: \newline
    When the I2S PCM-to-PDM converter is enabled, PDM clock frequency is equal to BCK frequency. The relation of sampling frequency ($f_{\mathrm{Sampling}}$, the sampling frequency of PCM data) and BCK frequency (the sampling frequency of PDM data) is as follows:
    $$f_{\mathrm{Sampling}} = \frac{f_{\mathrm{BCK}}}{\mathrm{OSR}}$$
    
    Upsampling rate (OSR) is related to \hyperref[fielddesc:I2STXPDMSINCOSR2]{I2S\_TX\_PDM\_SINC\_OSR2} as follows:
    $$\mathrm{OSR} = \mathrm{\hyperref[fielddesc:I2STXPDMSINCOSR2]{I2S\_TX\_PDM\_SINC\_OSR2}} \times 64$$
    
    Sampling frequency $f_{\mathrm{Sampling}}$ is related to \hyperref[fielddesc:I2STXPDMFS]{I2S\_TX\_PDM\_FS} as follows:
    $$f_{\mathrm{Sampling}} = \mathrm{\hyperref[fielddesc:I2STXPDMFS]{I2S\_TX\_PDM\_FS}} \times 100$$
    
    Configure the registers according to the needed sampling frequency, upsampling rate, and PDM clock frequency.
\end{itemize}

\textbf{PDM Channel Configuration Example}

In this example of I2S, the register configuration is as follows.

\begin{itemize}
    \item \hyperref[fielddesc:I2SPCM2PDMCONVEN]{I2S\_PCM2PDM\_CONV\_EN} = 0, i.e., transmit the raw PDM data.
    \item \hyperref[fielddesc:I2STXMONO]{I2S\_TX\_MONO} = 0, i.e., data is fetched from memory via GDMA in both the high and low levels of WS.
    \item \hyperref[fielddesc:I2STXCHANMOD]{I2S\_TX\_CHAN\_MOD} = 2, i.e., mono mode is selected.
    \item \hyperref[fielddesc:I2STXWSIDLEPOL]{I2S\_TX\_WS\_IDLE\_POL} = 1, i.e., both the left and right channels transmit the left channel data, and the right channel data will be discarded.
\end{itemize}

Once the configuration is done, assume that the data in memory \textbf{after data format control} is:

\begin{table}[H]
    \centering
    \caption{Data Arrangement After Data Format Control in PDM Channel Configuration Example}
    \label{tab:i2s-pdm-chan-example-data}
    \begin{tabular}{|c|c|c|c|c|c|c|}
    \hline
    Left & Right & Left & Right & ... & Left & Right \\ \hline
    \end{tabular}
\end{table}
\vspace{-0.5em}
\begin{tiplisting}
\vspace{-2em}
\begin{enumerate}
    \item The data above refers to the processed data after data format control instead of the original data.
    \item ``Left" and ``Right" represent channel data, and their bit widths are the channel valid data width. See Section \ref{sec:i2s-data-control}.
\end{enumerate}
\end{tiplisting}
\vspace{-1em}
Then the channel data is transmitted after channel mode control as follows.

\begin{figure}[H]
    \centering
    \includegraphics[width=0.6\textwidth]{\modulefiles/figures/i2s_chan_mode_20220210.pdf}
    \caption{PDM Channel Control Example}
    \label{fig:i2s-pdm-chan-control-example}
\end{figure}

\subsubsection{Synchronous Counter}

\chipname{} I2S TX unit contains synchronous counters to reflect the current I2S data transmission status. The counters can be used to detect I2S transmit data errors or to synchronize data when multiple devices are using I2S  for data transmission.

\chipname{} I2S contains two types of synchronous counters:

\begin{itemize}
    \item FIFO counter: \hyperref[fielddesc:I2STXFIFOCNT]{I2S\_TX\_FIFO\_CNT}
    \begin{itemize}
        \item Each time the I2S TX FIFO reads data, the counter is incremented by 1.
        \item The counter can be reset by setting \hyperref[fielddesc:I2STXFIFOCNTRST]{I2S\_TX\_FIFO\_CNT\_RST}.
        \item When ETM task I2S\regindex{n}\_TASK\_SYNC\_CHECK occurs, the counter is reset to zero.
        \item The counter is 31-bit wide and is automatically reset to zero on overflow.
    \end{itemize}
    \item BCK counter: \hyperref[fielddesc:I2STXBCKCNT]{I2S\_TX\_BCK\_CNT}
    \begin{itemize}
        \item Each time I2S TX transmits a BCK cycle, the counter is incremented by 1.
        \item The counter can be reset by setting \hyperref[fielddesc:I2STXBCKCNTRST]{I2S\_TX\_BCK\_CNT\_RST}.
        \item When ETM task I2S\regindex{n}\_TASK\_SYNC\_CHECK occurs, the counter is reset to zero.
        \item The counter is 31-bit wide and is automatically reset to zero on overflow.
    \end{itemize}
\end{itemize}

\vspace{-0.5em}
\begin{tiplisting}
\vspace{-2em}
\begin{itemize}
    \item Each time the I2S TX FIFO transmits data, a channel data is sent. Therefore, when I2S transmits the data correctly, \hyperref[fielddesc:I2STXBCKCNT]{I2S\_TX\_BCK\_CNT} = \hyperref[fielddesc:I2STXFIFOCNT]{I2S\_TX\_FIFO\_CNT} $\times$ the bit width of channel TX data (see Section \ref{sec:i2s-tdm-chan-bits}).
    \item When the I2S TX unit transmits the data, there is some delay between fetching data from the FIFO and transmitting it to the line, so the value of \hyperref[fielddesc:I2STXBCKCNT]{I2S\_TX\_BCK\_CNT} might be a little different from that of \hyperref[fielddesc:I2STXFIFOCNT]{I2S\_TX\_FIFO\_CNT}.
\end{itemize}
\end{tiplisting}
\vspace{-1em}

\textbf{Application Scenarios}
\begin{enumerate}
    \item I2S TX synchronization for multiple devices \\
    When multiple devices transmit data through the I2S TX at the same time, desynchronization may occur after a long period of time due to small frequency differences in the clock sources. In this case, read the \hyperref[fielddesc:I2STXFIFOCNT]{I2S\_TX\_FIFO\_CNT} value of each device to determine the location of the audio data being sent to perform synchronization.
    \item I2S TX stability check\\
    When the device is in an application scenario with high GDMA usage, there may be data missing in I2S TX occasionally due to multiplexing arbitration. In this case, check the ratio of \hyperref[fielddesc:I2STXBCKCNT]{I2S\_TX\_BCK\_CNT} to \hyperref[fielddesc:I2STXFIFOCNT]{I2S\_TX\_FIFO\_CNT} to confirm the stability of the device when running I2S TX.
\end{enumerate}

\paragraph{Hardware Automatic Synchronization}

The \chipname{} I2S\regindex{n} module supports hardware automatic synchronization, which works with ETM task I2S\regindex{n}\_TASK\_SYNC\_CHECK to calibrate I2S TX data automatically:

\begin{itemize}
    \item When I2S\regindex{n}\_TASK\_SYNC\_CHECK is executed, I2S reads \hyperref[fielddesc:I2STXFIFOCNT]{I2S\_TX\_FIFO\_CNT}, calculates its signed difference from the ideal value, and accumulates the result in \hyperref[fielddesc:I2STXCNTDIFF]{I2S\_TX\_CNT\_DIFF}:
    \[
        \hyperref[fielddesc:I2STXCNTDIFF]{I2S\_TX\_CNT\_DIFF} \mathrel{+}= \hyperref[fielddesc:I2STXFIFOCNT]{I2S\_TX\_FIFO\_CNT} - \hyperref[fielddesc:I2STXIDEALCNT]{I2S\_TX\_IDEAL\_CNT}
    \]
    The FIFO counter is then cleared and starts counting again. Thus, \hyperref[fielddesc:I2STXCNTDIFF]{I2S\_TX\_CNT\_DIFF} is the accumulated deviation since it was last cleared, rather than the deviation from one task execution. Set \hyperref[fielddesc:I2STXCNTDIFFRST]{I2S\_TX\_CNT\_DIFF\_RST} to clear it.
    \item If the absolute value of \hyperref[fielddesc:I2STXCNTDIFF]{I2S\_TX\_CNT\_DIFF} is greater than \hyperref[fielddesc:I2STXCNTDIFFHWTHRES]{I2S\_TX\_CNT\_DIFF\_HW\_THRES} and less than or equal to \hyperref[fielddesc:I2STXCNTDIFFSWTHRES]{I2S\_TX\_CNT\_DIFF\_SW\_THRES}, hardware corrects the deviation:
    \begin{itemize}
        \item A positive value means more data has been transmitted than the ideal value. Hardware supplements data as configured by \hyperref[fielddesc:I2STXHWSYNCSUPPLMODE]{I2S\_TX\_HW\_SYNC\_SUPPL\_MODE}; supplemented data is not fetched from the TX FIFO.
        \item A negative value means less data has been transmitted than the ideal value. Hardware additionally reads and discards data from the TX FIFO.
    \end{itemize}
    Each supplemented or discarded unit reduces the deviation by one until it reaches zero.
\end{itemize}

\begin{tiplisting}
\hyperref[fielddesc:I2STXCNTDIFF]{I2S\_TX\_CNT\_DIFF} is a 31-bit read-only field containing a signed difference in 31-bit two's-complement form; bit[30] is the sign bit:
\begin{itemize}
    \item bit[30] = 0: The difference is nonnegative and equals the read value.
    \item bit[30] = 1: The difference is negative and equals the read value minus $2^{31}$.
\end{itemize}
The representable range is $-2^{30}$ to $2^{30}-1$. Hardware uses the absolute value according to this interpretation when comparing thresholds; both threshold fields are unsigned.
\end{tiplisting}

The following configuration registers are used:

\begin{itemize}
    \item \hyperref[fielddesc:I2STXHWSYNCEN]{I2S\_TX\_HW\_SYNC\_EN}: Set to 1 to enable hardware automatic synchronization.
    \item \hyperref[fielddesc:I2STXIDEALCNT]{I2S\_TX\_IDEAL\_CNT}: Configures the ideal \hyperref[fielddesc:I2STXFIFOCNT]{I2S\_TX\_FIFO\_CNT} value when I2S\regindex{n}\_TASK\_SYNC\_CHECK is executed.
    \item \hyperref[fielddesc:I2STXCNTDIFFRST]{I2S\_TX\_CNT\_DIFF\_RST}: Set to clear the accumulated \hyperref[fielddesc:I2STXCNTDIFF]{I2S\_TX\_CNT\_DIFF}.
    \item \hyperref[fielddesc:I2STXCNTDIFFHWTHRES]{I2S\_TX\_CNT\_DIFF\_HW\_THRES}: Configures the lower deviation threshold for hardware automatic synchronization.
    \item \hyperref[fielddesc:I2STXCNTDIFFSWTHRES]{I2S\_TX\_CNT\_DIFF\_SW\_THRES}: Configures the upper deviation threshold. If the absolute deviation exceeds this value, an interrupt is triggered for software handling.

\begin{tiplisting}
\hyperref[fielddesc:I2STXCNTDIFFHWTHRES]{I2S\_TX\_CNT\_DIFF\_HW\_THRES} and \hyperref[fielddesc:I2STXCNTDIFFSWTHRES]{I2S\_TX\_CNT\_DIFF\_SW\_THRES} together determine whether synchronization is handled automatically by hardware or by software:
    \begin{itemize}
        \item If the \textbf{absolute value} of \hyperref[fielddesc:I2STXCNTDIFF]{I2S\_TX\_CNT\_DIFF} is less than or equal to \hyperref[fielddesc:I2STXCNTDIFFHWTHRES]{I2S\_TX\_CNT\_DIFF\_HW\_THRES}, no synchronization action is taken.
        \item If the absolute value of \hyperref[fielddesc:I2STXCNTDIFF]{I2S\_TX\_CNT\_DIFF} is greater than \hyperref[fielddesc:I2STXCNTDIFFHWTHRES]{I2S\_TX\_CNT\_DIFF\_HW\_THRES} and less than or equal to \hyperref[fielddesc:I2STXCNTDIFFSWTHRES]{I2S\_TX\_CNT\_DIFF\_SW\_THRES}, hardware handles synchronization automatically.
        \item If the absolute value of \hyperref[fielddesc:I2STXCNTDIFF]{I2S\_TX\_CNT\_DIFF} exceeds \hyperref[fielddesc:I2STXCNTDIFFSWTHRES]{I2S\_TX\_CNT\_DIFF\_SW\_THRES}, hardware triggers an interrupt and hands over synchronization to software.
    \end{itemize}
\end{tiplisting}

    \item \hyperref[fielddesc:I2STXHWSYNCSUPPLMODE]{I2S\_TX\_HW\_SYNC\_SUPPL\_MODE}:
    \begin{itemize}
        \item 0: Automatically supplement the last transmitted data.
        \item 1: Automatically supplement the data specified in \hyperref[fielddesc:I2STXHWSYNCSUPPLDATA]{I2S\_TX\_HW\_SYNC\_SUPPL\_DATA}.
    \end{itemize}
\end{itemize}

\subsection{Receiving Data}\label{sec:i2s-rx-data}

In RX mode, I2S first reads data from the peripheral interface and then stores the data in memory via GDMA according to the configured channel mode and data mode.

\subsubsection{Channel Mode Control}

\chipname{} I2S supports both TDM RX mode and PDM RX mode. Set \hyperref[fielddesc:I2SRXTDMEN]{I2S\_RX\_TDM\_EN} to enable TDM RX mode, or set \hyperref[fielddesc:I2SRXPDMEN]{I2S\_RX\_PDM\_EN} to enable PDM RX mode. PDM RX mode is available only on I2S\regindex{0}.

\begin{tiplisting}
\hyperref[fielddesc:I2SRXTDMEN]{I2S\_RX\_TDM\_EN} and \hyperref[fielddesc:I2SRXPDMEN]{I2S\_RX\_PDM\_EN} must not be cleared or set simultaneously.
\end{tiplisting}
\vspace{-1em}

\paragraph{TDM RX Mode}

In TDM RX mode, I2S supports up to 16 channels to input data. The total number of RX channels in use is controlled by \hyperref[fielddesc:I2SRXTDMTOTCHANNUM]{I2S\_RX\_TDM\_TOT\_CHAN\_NUM}. For example, if \hyperref[fielddesc:I2SRXTDMTOTCHANNUM]{I2S\_RX\_TDM\_TOT\_CHAN\_NUM} is set to 5, channels 0--5 will be used to receive data.

For enabled receive channels, the channel-enable fields are \hyperref[fielddesc:I2SRXTDMPDMCHAN0EN]{I2S\_RX\_TDM\_PDM\_CHAN{\textcolor{red}{n}}\_EN} for channels 0--7 and \hyperref[fielddesc:I2SRXTDMCHAN8EN]{I2S\_RX\_TDM\_CHAN{\textcolor{red}{n}}\_EN} for channels 8--15. If the corresponding field is:
\begin{itemize}
    \item 1: The channel data is valid and will be stored into RX FIFO.
    \item 0: The channel data is invalid and will not be stored into RX FIFO.
\end{itemize}
In TDM master mode, WS signal is controlled by \hyperref[fielddesc:I2SRXWSIDLEPOL]{I2S\_RX\_WS\_IDLE\_POL} and \hyperref[fielddesc:I2SRXTDMWSWIDTH]{I2S\_RX\_TDM\_WS\_WIDTH}.
\begin{itemize}
    \item \hyperref[fielddesc:I2SRXWSIDLEPOL]{I2S\_RX\_WS\_IDLE\_POL}: The default level of WS signal.
    \item \hyperref[fielddesc:I2SRXTDMWSWIDTH]{I2S\_RX\_TDM\_WS\_WIDTH}: Configures the idle duration during which WS remains at its default level during reception. Idle duration = (\hyperref[fielddesc:I2SRXTDMWSWIDTH]{I2S\_RX\_TDM\_WS\_WIDTH} + 1) × T\_BCK.
\end{itemize}

(\hyperref[fielddesc:I2SRXHALFSAMPLEBITS]{I2S\_RX\_HALF\_SAMPLE\_BITS} + 1) × 2 is equal to the BCK cycles in one WS period.

\paragraph{PDM RX Mode} 

PDM RX mode is available only on I2S\regindex{0}. In this mode, I2S\regindex{0} can receive raw PDM data or convert received PDM data to PCM data.

In PDM RX master mode, the default level of the WS signal is controlled by \hyperref[fielddesc:I2SRXWSIDLEPOL]{I2S\_RX\_WS\_IDLE\_POL}. WS frequency is half of the BCK frequency. The configuration of the BCK signal is similar to that of WS signal as described in Section \ref{sec:i2s-tx-rx-clk}. Note, in PDM RX mode, the value of \hyperref[fielddesc:I2SRXHALFSAMPLEBITS]{I2S\_RX\_HALF\_SAMPLE\_BITS} must be same as that of \hyperref[fielddesc:I2SRXBITSMOD]{I2S\_RX\_BITS\_MOD}.

When the PDM-to-PCM converter is enabled, I2S\regindex{0} converts received PDM data to PCM data before subsequent processing. Set \hyperref[fielddesc:I2SRXPDM2PCMEN]{I2S\_RX\_PDM2PCM\_EN} to enable the converter.
\begin{itemize}
    \item Configure the sampling frequency and downsampling rate. In master mode, the PDM clock frequency is the BCK frequency; in slave mode, it is provided externally:
    $$f_{\mathrm{Sampling}} = \frac{f_{\mathrm{PDM}}}{\mathrm{64 \times DSR}}$$
    The downsampling rate is:
    $$\mathrm{DSR} = 8 \times 2^{\mathrm{\hyperref[fielddesc:I2SRXPDMSINCDSR16EN]{I2S\_RX\_PDM\_SINC\_DSR\_16\_EN}}}$$
    \item Configure valid channels. The converter supports up to four PDM data lines and eight input channels. Enable channels with \hyperref[fielddesc:I2SRXTDMPDMCHAN0EN]{I2S\_RX\_TDM\_PDM\_CHAN\regindex{n}\_EN} (\regindex{n}: 0--7). Each data line carries two channels in one WS period: the even channel is sampled on the rising edge of WS and the odd channel on the falling edge.
\end{itemize}

\begin{table}[H]
    \centering
    \caption{PDM-to-PCM Converter Input Channel Mapping}
    \label{tab:i2s-pdm2pcm-chan-map}
    \begin{tabular}{|p{3.6cm}|R{4.2cm}|R{4.2cm}|}
    \hline
    \rowcolor{lightgray}
    \textbf{Data Line} & \textbf{Channel Sampled on WS Rising Edge} & \textbf{Channel Sampled on WS Falling Edge}\\ \hline
    I2S\regindex{0}I\_SD\_in  & Channel 0 & Channel 1 \\\hline
    I2S\regindex{0}I\_SD1\_in & Channel 2 & Channel 3 \\\hline
    I2S\regindex{0}I\_SD2\_in & Channel 4 & Channel 5 \\\hline
    I2S\regindex{0}I\_SD3\_in & Channel 6 & Channel 7 \\\hline
    \end{tabular}
\end{table}

\subsubsection{Data Format Control}

The data format of I2S is controlled in the following phases:

\begin{itemize}
    \item Phase I: Serial input data is converted into the data to be saved to RX FIFO.
    \item Phase II: The data is read from RX FIFO and converted according to the input data mode.
\end{itemize}

\paragraph{Channel Data Bit Order} 

The channel data will be stored as the data to be input in order from high to low. The data bit order in each channel is controlled by  \hyperref[fielddesc:I2SRXBITORDER]{I2S\_RX\_BIT\_ORDER}:
\begin{itemize}
\item 0: The bit order of the data to be input is not reversed.
\item 1: The bit order of the data to be input is reversed.
\end{itemize}
At this point, the first phase of data format control is completed. The data to be input after bit order control is stored in the RX FIFO.

\paragraph{Channel Storage Data Width}

The storage data width in each channel is controlled by \hyperref[fielddesc:I2SRXBITSMOD]{I2S\_RX\_BITS\_MOD} and \hyperref[fielddesc:I2SRX24FILLEN]{I2S\_RX\_24\_FILL\_EN}. See the table below.
\begin{table}[H]
    \centering
    \caption{Channel Storage Data Width}
    \label{tab:i2s-chan-storage-data-width}
    \begin{tabular}{|p{5cm}|R{3.3cm}|R{4cm}|}
    \hline
    \rowcolor{lightgray}
    \textbf{Channel Storage Data Width} & \textbf{\hyperref[fielddesc:I2SRXBITSMOD]{I2S\_RX\_BITS\_MOD}} & \textbf{\hyperref[fielddesc:I2SRX24FILLEN]{I2S\_RX\_24\_FILL\_EN}}  \\ \hline
    \multirow{2}*{32} & 31 & x \\\cline{2-3}
                 & 23 & 1 \\ \hline
    24 & 23 & 0 \\\hline
    16    & 15 & x \\\hline
    8     & 7 & x \\\hline
    \end{tabular}
\end{table}

\paragraph{Channel RX Data Width} 

The RX data width in each channel is determined by \hyperref[fielddesc:I2SRXTDMCHANBITS]{I2S\_RX\_TDM\_CHAN\_BITS}.
\begin{itemize}
    \item If the storage data width in each channel is smaller than the received (RX) data width, then only the bits within the storage data width is saved into memory. Configure \hyperref[fielddesc:I2SRXLEFTALIGN]{I2S\_RX\_LEFT\_ALIGN} to:
        \begin{itemize}
        \item 0: Only the lower bits of the received data within the storage data width is stored to memory.
        \item 1: Only the higher bits of the received data within the storage data width is stored to memory.
        \end{itemize}
    \item If the received data width is smaller than the storage data width in each channel, the higher bits of the received data will be filled with zeros and then the data is saved to memory.
\end{itemize}

\paragraph{Channel Storage Data Endianness} 

The received data is then converted into storage data (to be stored to memory) after some processing, such as discarding extra bits or filling zeros in missing bits. The endian of the storage data is controlled by \hyperref[fielddesc:I2SRXBIGENDIAN]{I2S\_RX\_BIG\_ENDIAN} under various data width. See the table below.

\begin{table}[H]
    \centering
    \caption{Channel Storage Data Endian}
    \label{tab:i2s-chan-storage-data-endian}
    \begin{tabular}{|p{4.8cm}|R{2.5cm}|R{4.2cm}|R{3.5cm}|}
    \hline
    \rowcolor{lightgray}
    \textbf{Channel Storage Data Width} & \textbf{Original Data} & \textbf{Endian of Processed Data} & \textbf{\hyperref[fielddesc:I2SRXBIGENDIAN]{I2S\_RX\_BIG\_ENDIAN}} \\ \hline
 \multirow{2}*{32} & \multirow{2}*{\{B3, B2, B1, B0\}} & \{B3, B2, B1, B0\} & 0 \\\cline{3-4}
                      && \{B0, B1, B2, B3\} & 1 \\\hline

    \multirow{2}*{24} & \multirow{2}*{\{B2, B1, B0\}} & \{B2, B1, B0\} & 0 \\\cline{3-4}
                      && \{B0, B1, B2\} & 1 \\\hline

    \multirow{2}*{16} & \multirow{2}*{\{B1, B0\}} & \{B1, B0\} & 0 \\\cline{3-4}
                      && \{B0, B1\} & 1 \\\hline

                8     & \{B0\} & \{B0\} & x \\\hline
    \end{tabular}
\end{table}

\paragraph{A-law/$\mu$-law Compression and Decompression} 

\chipname{} I2S compresses/decompresses the storage data in 32-bit by A-law or by $\mu$-law. By default, zeros are filled into high bits.

Configure \hyperref[fielddesc:I2SRXPCMBYPASS]{I2S\_RX\_PCM\_BYPASS}:
\begin{itemize}
\item 0: Compress or decompress the data
\item 1: Do not compress or decompress the data
\end{itemize}
Configure \hyperref[fielddesc:I2SRXPCMCONF]{I2S\_RX\_PCM\_CONF}:
\begin{itemize}
\item 0: Decompress the data using A-law
\item 1: Compress the data using A-law
\item 2: Decompress the data using $\mu$-law
\item 3: Compress the data using $\mu$-law
\end{itemize}

At this point, the data format control is completed. Data then is stored into memory via GDMA.

\subsubsection{RX Data Recombination}\label{sec:i2s-rx-recomb}

The I2S\regindex{n} RX unit outputs data to GDMA through DMA interfaces. By default, formatted RX data is output as one stream through DMA interface 0. Data from different receive channels is interleaved in memory and must be separated by software.

I2S\regindex{0} supports RX data recombination. When enabled, the RX stream is first separated by receive channel and then recombined, as configured by software, into up to four streams. These streams are output through four DMA interfaces and stored at independent memory locations. Therefore, different software modules can obtain the channel data they require without moving data in memory.

Set \hyperref[fielddesc:I2SRXRECOMBEN]{I2S\_RX\_RECOMB\_EN} to enable this feature. Recombination consists of channel extraction and channel combination.

\begin{tiplisting}
DMA interfaces 0--3 in this subsection are the four DMA interfaces provided by I2S\regindex{0}. Their numbers correspond to \regindex{n} in \hyperref[regdesc:I2SRXRECOMBDMACH0REG]{I2S\_RX\_RECOMB\_DMA\_CH\regindex{n}\_REG} and are unrelated to GDMA internal channel numbers. GDMA uses a shared-channel architecture; the mapping between its internal channels and peripheral DMA interfaces is configured in GDMA. See Chapter \ref{mod:dma} \textit{\nameref{mod:dma}}.
\end{tiplisting}

\paragraph{Channel Extraction}

The RX stream is arranged cyclically by receive channel in 32-bit units. The recombination unit distributes the units in order to the input-channel buffers: the first unit enters input channel 0, the second enters input channel 1, and so on. After the last input channel, distribution returns to input channel 0.

\hyperref[fielddesc:I2SRXRECOMBEXTCHNUM]{I2S\_RX\_RECOMB\_EXT\_CH\_NUM} configures the number of input channels minus one:
\begin{itemize}
    \item 0: One input channel
    \item 1: Two input channels
    \item 2: Three input channels
    \item 3: Four input channels
\end{itemize}

\begin{tiplisting}
\begin{itemize}
    \item The RX stream contains data only from enabled receive channels (see \hyperref[fielddesc:I2SRXTDMPDMCHAN0EN]{I2S\_RX\_TDM\_PDM\_CHAN\regindex{n}\_EN}). Therefore, configure \hyperref[fielddesc:I2SRXRECOMBEXTCHNUM]{I2S\_RX\_RECOMB\_EXT\_CH\_NUM} according to the actual number of enabled receive channels, not \hyperref[fielddesc:I2SRXTDMTOTCHANNUM]{I2S\_RX\_TDM\_TOT\_CHAN\_NUM}.
    \item Recombination supports at most four input channels. Do not use it if more than four receive channels are enabled.
    \item Data from an input channel that is not selected by any DMA interface is discarded.
\end{itemize}
\end{tiplisting}

\paragraph{Channel Combination}

Each DMA interface has one \hyperref[regdesc:I2SRXRECOMBDMACH0REG]{I2S\_RX\_RECOMB\_DMA\_CH\regindex{n}\_REG} register (\regindex{n}: 0--3), which configures the input channels and output order:
\begin{itemize}
    \item \hyperref[fielddesc:I2SRXRECOMBDMACH0VALID]{I2S\_RX\_RECOMB\_DMA\_CH\regindex{n}\_VALID}: Configures whether this DMA interface participates in recombination.
    \begin{itemize}
        \item 0: Does not participate and outputs no data
        \item 1: Participates
    \end{itemize}
    \item \hyperref[fielddesc:I2SRXRECOMBDMACH0STYLE]{I2S\_RX\_RECOMB\_DMA\_CH\regindex{n}\_STYLE}: Selects input channels for this DMA interface. This is a 4-bit bitmap, where bit \regindex{i} corresponds to input channel \regindex{i}. For example, configure 4'b1010 to select input channels 1 and 3.
    \item \hyperref[fielddesc:I2SRXRECOMBDMACH0ORDER]{I2S\_RX\_RECOMB\_DMA\_CH\regindex{n}\_ORDER}: Configures the output order. This 8-bit field contains four 2-bit subfields, starting at the least-significant bits, that specify the first, second, and subsequent input-channel numbers. For example, to output input channel 3, then channel 1, then channel 0, and then channel 2, configure 8'b1000\_0111 (bit[1:0] = 3, bit[3:2] = 1, bit[5:4] = 0, and bit[7:6] = 2).
    \item \hyperref[fielddesc:I2SRXRECOMBDMACH0EOFNUM]{I2S\_RX\_RECOMB\_DMA\_CH\regindex{n}\_EOF\_NUM}: Configures the amount of data received by this DMA interface in one transfer, in bytes.
\end{itemize}

Each DMA interface reads and outputs one unit from each selected input-channel buffer in the order specified by ORDER. After all channels selected by STYLE have been read, the interface starts the next round. The same input channel can be selected by multiple DMA interfaces.

\begin{tiplisting}
\begin{itemize}
\vspace{-2em}
    \item The first applicable ORDER subfields must contain each input channel selected by STYLE exactly once; remaining subfields are ignored. If ORDER specifies a channel not selected by STYLE, the recombination unit stops fetching data.
    \item After modifying any register above, set \hyperref[fielddesc:I2SRXRECOMBUPDATE]{I2S\_RX\_RECOMB\_UPDATE} to synchronize the configuration from the APB clock domain to the I2S RX clock domain. Hardware clears this bit automatically.
\end{itemize}
\end{tiplisting}

\paragraph{Length Control and Transfer Completion}

Each DMA interface counts its output bytes independently. When its count reaches \hyperref[fielddesc:I2SRXRECOMBDMACH0EOFNUM]{I2S\_RX\_RECOMB\_DMA\_CH\regindex{n}\_EOF\_NUM}, the interface sends an end-of-transfer indication to GDMA. Software can use the corresponding GDMA channel transfer-completion interrupt to determine that reception is complete.

DMA interface lengths and completion are independent. Set \hyperref[fielddesc:I2SRXRECOMBDMACH0VALID]{I2S\_RX\_RECOMB\_DMA\_CH\regindex{n}\_VALID} for each interface in use and leave unused interfaces cleared.

\paragraph{Configuration Example}

Assume RX has four enabled receive channels. DMA interface 0 selects input channel 0; DMA interface 1 selects input channel 1; DMA interface 2 selects input channel 3 and then input channel 2; and DMA interface 3 selects input channel 0 and then input channel 1. Table \ref{tab:i2s-rx-recomb-example} shows the configuration.

\begin{table}[H]
    \centering
    \caption{RX Data Recombination Configuration Example}
    \label{tab:i2s-rx-recomb-example}
    \begin{tabular}{|p{6cm}|R{2.0cm}|R{2.6cm}|R{3.0cm}|}
    \hline
    \rowcolor{lightgray}
    \textbf{Register Group} & \textbf{\hyperref[fielddesc:I2SRXRECOMBDMACH0VALID]{VALID}} & \textbf{\hyperref[fielddesc:I2SRXRECOMBDMACH0STYLE]{STYLE}} & \textbf{\hyperref[fielddesc:I2SRXRECOMBDMACH0ORDER]{ORDER}}\\ \hline
    \hyperref[regdesc:I2SRXRECOMBDMACH0REG]{I2S\_RX\_RECOMB\_DMA\_CH0\_REG} & 1 & 4'b0001 & 8'b0000\_0000 \\\hline
    \hyperref[regdesc:I2SRXRECOMBDMACH1REG]{I2S\_RX\_RECOMB\_DMA\_CH1\_REG} & 1 & 4'b0010 & 8'b0000\_0001 \\\hline
    \hyperref[regdesc:I2SRXRECOMBDMACH2REG]{I2S\_RX\_RECOMB\_DMA\_CH2\_REG} & 1 & 4'b1100 & 8'b0000\_1011 \\\hline
    \hyperref[regdesc:I2SRXRECOMBDMACH3REG]{I2S\_RX\_RECOMB\_DMA\_CH3\_REG} & 1 & 4'b0011 & 8'b0000\_0100 \\\hline
    \end{tabular}
\end{table}

Configure \hyperref[fielddesc:I2SRXRECOMBEXTCHNUM]{I2S\_RX\_RECOMB\_EXT\_CH\_NUM} to 3 for four input channels. For DMA interface 2, ORDER has bit[1:0] = 2'd3 and bit[3:2] = 2'd2. For DMA interface 3, ORDER has bit[1:0] = 2'd0 and bit[3:2] = 2'd1. Configure each interface's \hyperref[fielddesc:I2SRXRECOMBDMACH0EOFNUM]{I2S\_RX\_RECOMB\_DMA\_CH\regindex{n}\_EOF\_NUM} for the required number of bytes.

After configuration, set \hyperref[fielddesc:I2SRXRECOMBUPDATE]{I2S\_RX\_RECOMB\_UPDATE}, and then start reception as described in Section \ref{sec:i2s-configure-i2s-rx-mode} \textit{\nameref{sec:i2s-configure-i2s-rx-mode}}.

\subsection{Clock}\label{sec:i2s-tx-rx-clk}

I2S\regindex{n}\_TX/RX\_CLK is the main clock of the I2S\regindex{n} TX/RX unit and can be divided from:
\begin{itemize}
    \item XTAL\_CLK
    \item APLL\_CLK
    \item RC\_FAST\_CLK
    \item External input clock I2S\regindex{n}\_MCLK\_in
\end{itemize}

The serial BCK clock is divided from I2S\regindex{n}\_TX/RX\_CLK, as shown in Figure \ref{fig:i2s-clk}. Select and enable the I2S TX/RX clock source with \hyperref[fielddesc:HPSYSCLKRSTREGI2S0RXCLKSRCSEL]{HP\_SYS\_CLKRST\_REG\_I2S\regindex{n}\_RX/TX\_CLK\_SRC\_SEL} and \hyperref[fielddesc:HPSYSCLKRSTREGI2S0RXCLKEN]{HP\_SYS\_CLKRST\_REG\_I2S\regindex{n}\_RX/TX\_CLK\_EN}. Configure the I2S\regindex{n} MCLK output clock with \hyperref[fielddesc:HPSYSCLKRSTREGI2S0MSTCLKSEL]{HP\_SYS\_CLKRST\_REG\_I2S\regindex{n}\_MST\_CLK\_SEL}.

\begin{figure}[H]
    \centering
    \includegraphics[width=0.9\textwidth]{\modulefiles/figures/i2s_clk_20260818}
    \caption{I2S Clock Divider}
    \label{fig:i2s-clk}
\end{figure}

The relation between $f_{\mathrm{I2S\regindex{n}\_TX/RX\_CLK}}$ and the selected source frequency $f_{\mathrm{I2S\regindex{n}\_CLK\_S}}$ is:
$$f_{\mathrm{I2S\regindex{n}\_TX/RX\_CLK}} = \frac{f_{\mathrm{I2S\regindex{n}\_CLK\_S}}}{N + \frac{b}{a}}$$

Here, $2 \leq N \leq 256$. N maps to \hyperref[fielddesc:HPSYSCLKRSTREGI2S0RXDIVN]{HP\_SYS\_CLKRST\_REG\_I2S\regindex{n}\_RX/TX\_DIV\_N} as follows:
\begin{itemize}
    \item Field value 0: N = 256
    \item Field value 1: N = 2
    \item Other field values: N equals the field value
\end{itemize}

The fractional coefficients a and b uniquely determine x, y, z, and yn1:
\begin{itemize}
    \item If $b \leq \frac{a}{2}$, $\mathrm{yn1} = 0$, $x = \operatorname{floor}(\frac{a}{b}) - 1$, $y = a \bmod b$, and $z = b$.
    \item If $b > \frac{a}{2}$, $\mathrm{yn1} = 1$, $x = \operatorname{floor}(\frac{a}{a - b}) - 1$, $y = a \bmod (a - b)$, and $z = a - b$.
\end{itemize}

Configure x, y, z, and yn1 with \hyperref[fielddesc:HPSYSCLKRSTREGI2S0RXDIVX]{HP\_SYS\_CLKRST\_REG\_I2S\regindex{n}\_RX/TX\_DIV\_X}, \hyperref[fielddesc:HPSYSCLKRSTREGI2S0RXDIVY]{HP\_SYS\_CLKRST\_REG\_I2S\regindex{n}\_RX/TX\_DIV\_Y}, \hyperref[fielddesc:HPSYSCLKRSTREGI2S0RXDIVZ]{HP\_SYS\_CLKRST\_REG\_I2S\regindex{n}\_RX/TX\_DIV\_Z}, and \hyperref[fielddesc:HPSYSCLKRSTREGI2S0RXDIVYN1]{HP\_SYS\_CLKRST\_REG\_I2S\regindex{n}\_RX/TX\_DIV\_YN1}. For integer division, clear DIV\_X and DIV\_Z and set DIV\_Y to 1.

\begin{tiplisting}
\begin{itemize}
    \item The maximum $f_{\mathrm{I2S\regindex{n}\_TX/RX\_CLK}}$ is 62.5 MHz.
    \item Fractional division may introduce clock jitter.
\end{itemize}
\end{tiplisting}

In master TX mode:
$$f_{\mathrm{I2S\regindex{n}O\_BCK\_out}} = \frac{f_{\mathrm{I2S\regindex{n}\_TX\_CLK}}}{\mathrm{\hyperref[fielddesc:I2STXBCKDIVNUM]{I2S\_TX\_BCK\_DIV\_NUM}} + 1}$$
\begin{tiplisting}
\hyperref[fielddesc:I2STXBCKDIVNUM]{I2S\_TX\_BCK\_DIV\_NUM} must not be configured as 1.
\end{tiplisting}

In master RX mode:
$$f_{\mathrm{I2S\regindex{n}I\_BCK\_out}} = \frac{f_{\mathrm{I2S\regindex{n}\_RX\_CLK}}}{\mathrm{\hyperref[fielddesc:I2SRXBCKDIVNUM]{I2S\_RX\_BCK\_DIV\_NUM}} + 1}$$
\begin{tiplisting}
\begin{itemize}
    \item \hyperref[fielddesc:I2SRXBCKDIVNUM]{I2S\_RX\_BCK\_DIV\_NUM} must not be configured as 1.
    \item In slave mode, $f_{\mathrm{I2S\regindex{n}\_TX/RX\_CLK}} \geq 8 \times f_{\mathrm{BCK}}$. The unit can also output I2S\_MCLK\_out as the master clock for an external device.
\end{itemize}
\end{tiplisting}

\subsection{Reset}\label{sec:i2s-reset}

Reset the I2S units and FIFOs with the following fields:
\begin{itemize}
    \item I2S\regindex{n} APB clock and reset: Enable the APB clock with \hyperref[fielddesc:HPSYSCLKRSTREGI2S0APBCLKEN]{HP\_SYS\_CLKRST\_REG\_I2S\regindex{n}\_APB\_CLK\_EN}, trigger APB reset with \hyperref[fielddesc:HPSYSCLKRSTREGI2S0APBRSTEN]{HP\_SYS\_CLKRST\_REG\_I2S\regindex{n}\_APB\_RST\_EN}, and configure whether reset is masked with \hyperref[fielddesc:HPSYSCLKRSTREGI2S0FORCENORST]{HP\_SYS\_CLKRST\_REG\_I2S\regindex{n}\_FORCE\_NORST}.
    \item I2S\regindex{n} TX/RX units: Reset with \hyperref[fielddesc:I2STXRESET]{I2S\_TX\_RESET} and \hyperref[fielddesc:I2SRXRESET]{I2S\_RX\_RESET}.
    \item I2S\regindex{n} TX/RX FIFOs: Reset with \hyperref[fielddesc:I2STXFIFORESET]{I2S\_TX\_FIFO\_RESET} and \hyperref[fielddesc:I2SRXFIFORESET]{I2S\_RX\_FIFO\_RESET}.
\end{itemize}
\begin{tiplisting}
Configure the I2S module clock before resetting the module or FIFO.
\end{tiplisting}

\section{Event Task Matrix Feature}\label{sec:i2s-etm-func}

\chipname{} I2S supports the Event Task Matrix (ETM) function, which allows I2S's ETM tasks to be triggered by any peripherals' ETM events, or I2S's ETM events to trigger any peripherals' ETM tasks. This section introduces the ETM tasks and events related to I2S. For more information, please refer to Chapter \ref{mod:evntaskmatrix} \nameref{mod:evntaskmatrix}.

I2S can receive the following ETM tasks:

\begin{itemize}
    \item \label{task:I2STASKSTARTTX}I2S\regindex{n}\_TASK\_START\_TX: Starts I2S TX data transfer.
    \item \label{task:I2STASKSTARTRX}I2S\regindex{n}\_TASK\_START\_RX: Starts I2S RX data transfer.
    \item \label{task:I2STASKSTOPTX}I2S\regindex{n}\_TASK\_STOP\_TX: Stops I2S TX data transfer.
    \item \label{task:I2STASKSTOPRX}I2S\regindex{n}\_TASK\_STOP\_RX: Stops I2S RX data transfer.
    \item \label{task:I2STASKSYNCCHECK}I2S\regindex{n}\_TASK\_SYNC\_CHECK: Triggers the I2S TX synchronization-counter function.
\end{itemize}

I2S can generate the following ETM events:

\begin{itemize}
    \item \label{evt:I2SEVTTXDONE}I2S\regindex{n}\_EVT\_TX\_DONE: Indicates that all data in the TX FIFO has been transmitted.
    \item \label{evt:I2SEVTRXDONE}I2S\regindex{n}\_EVT\_RX\_DONE: Its trigger condition depends on \hyperref[fielddesc:I2SRXSTOPMODE]{I2S\_RX\_STOP\_MODE}:
    \begin{itemize}
        \item 0: Will not be triggered.
        \item 1: Triggered when the number of bytes received reaches the length configured in \hyperref[regdesc:I2SRXEOFNUMREG]{I2S\_RXEOF\_NUM\_REG}.
        \item 2: Uses the same trigger condition as \hyperref[int:I2SRXDONEINT]{I2S\regindex{n}\_RX\_DONE\_INT}, configured by \hyperref[fielddesc:I2SRXDONEMODE]{I2S\_RX\_DONE\_MODE}.
    \end{itemize}
    \item \label{evt:I2SEVTXWORDSSENT}I2S\regindex{n}\_EVT\_X\_WORDS\_SENT: Indicates that the number of words transmitted is equal to or greater than \hyperref[fielddesc:I2SETMTXSENDWORDNUM]{I2S\_ETM\_TX\_SEND\_WORD\_NUM}.
    \item \label{evt:I2SEVTXWORDSRECEIVED}I2S\regindex{n}\_EVT\_X\_WORDS\_RECEIVED: Indicates that the number of words received is equal to or greater than \hyperref[fielddesc:I2SETMRXRECEIVEWORDNUM]{I2S\_ETM\_RX\_RECEIVE\_WORD\_NUM}.
\end{itemize}

For example, I2S\regindex{n}\_EVT\_X\_WORDS\_SENT can trigger I2S\regindex{n}\_TASK\_STOP\_TX to stop I2S\regindex{n} through ETM.


%%%%%%%%%%%%%%%%%%%%%
%%% Interrupts %%%
%%%%%%%%%%%%%%%%%%%%%

\section{Interrupts} \label{sec:i2s-interrupts}

\chipname{} I2S\regindex{n} can generate the I2S\_INTR interrupt signal and send it to the \hyperref[mod:intmtrx]{Interrupt Matrix}. This interrupt signal is generated by the following I2S\regindex{n} internal interrupt sources.

\begin{itemize}
\item \label{int:I2STXHUNGINT} I2S\regindex{n}\_TX\_HUNG\_INT: Triggered when data transmission times out. For example, if I2S is configured as TX slave mode but the master does not provide BCK or WS signals for a long time, this interrupt is triggered. Timeout configuration is in \hyperref[regdesc:I2SLCHUNGCONFREG]{I2S\_LC\_HUNG\_CONF\_REG}.

\item \label{int:I2SRXHUNGINT} I2S\regindex{n}\_RX\_HUNG\_INT: Triggered when data reception times out. For example, if I2S is configured as RX slave mode but the master does not transmit data for a long time, this interrupt is triggered. Timeout configuration is in \hyperref[regdesc:I2SLCHUNGCONFREG]{I2S\_LC\_HUNG\_CONF\_REG}.
\item \label{int:I2STXDONEINT} I2S\regindex{n}\_TX\_DONE\_INT: Triggered when data transmission is complete, that is, all data in the I2S TX FIFO has been sent.
\item \label{int:I2SRXDONEINT} I2S\regindex{n}\_RX\_DONE\_INT: Triggered according to \hyperref[fielddesc:I2SRXDONEMODE]{I2S\_RX\_DONE\_MODE}:
    \begin{itemize}
        \item 0: Triggered when the GDMA RX FIFO is full.
        \item 1: Triggered when the number of received bytes reaches the length configured in \hyperref[regdesc:I2SRXEOFNUMREG]{I2S\_RXEOF\_NUM\_REG}.
    \end{itemize}
\item \label{int:I2STXSYNCINT}\hyperref[int:I2STXSYNCINT]{I2S\regindex{n}\_TX\_SYNC\_INT}: Triggered when the absolute value of \hyperref[fielddesc:I2STXCNTDIFF]{I2S\_TX\_CNT\_DIFF} exceeds \hyperref[fielddesc:I2STXCNTDIFFSWTHRES]{I2S\_TX\_CNT\_DIFF\_SW\_THRES}; software must handle synchronization.
\end{itemize}

\begin{tiplisting}
For definitions of \textit{interrupt}, \textit{interrupt signal}, \textit{interrupt source}, and their correlations, please refer to Chapter \ref{mod:intmtrx} \textit{\nameref{mod:intmtrx}} > Section \ref{sec:intmtx-int-terminology} \textit{\nameref{sec:intmtx-int-terminology}}.
\end{tiplisting}

Each interrupt source can be configured by a common set of registers that are described in Section \hyperref[interrupt-config-registers]{\textit{Interrupt Configuration Registers}}. The specific registers can be found in Section \ref{sec:i2s-reg-summ} \textit{\nameref{sec:i2s-reg-summ}}.

\section{Programming Procedures}
\label{sec:i2s-prog-procedures}
\subsection{Configure I2S as TX Mode}\label{sec:i2s-configure-i2s-tx-mode}
Follow the steps below to configure I2S as TX mode via software:

\begin{enumerate}
    \item Configure the clock as described in Section \ref{sec:i2s-tx-rx-clk}.
    \item Configure signal pins according to Table \ref{tab:i2s-signal-descr}.
    \item Select the mode needed by configuring \hyperref[fielddesc:I2STXSLAVEMOD]{I2S\_TX\_SLAVE\_MOD}.
    \begin{itemize}
        \item 0: Master TX mode
        \item 1: Slave TX mode
    \end{itemize}
    \item Set needed TX data mode and TX channel mode as described in Section \ref{sec:i2s-tx-data}, and then set \hyperref[fielddesc:I2STXUPDATE]{I2S\_TX\_UPDATE}.
    \item Reset TX unit and TX FIFO as described in Section \ref{sec:i2s-reset}.
    \item Enable corresponding interrupts. See Section \ref{sec:i2s-interrupts}.\label{other:i2s-item-tx-enable-corresponding-interrupts}
    \item Configure GDMA outlink. 
    \item Set \hyperref[fielddesc:I2STXSTOPEN]{I2S\_TX\_STOP\_EN} if needed. For more information, please refer to Section \ref{sec:i2s-master-slave-tx-mode}. 
    \item Start transmitting data:
    \begin{itemize}
        \item In master mode, wait till I2S slave gets ready, then set \hyperref[fielddesc:I2STXSTART]{I2S\_TX\_START} to start transmitting data.
        \item In slave mode, set \hyperref[fielddesc:I2STXSTART]{I2S\_TX\_START}. When the I2S master supplies BCK and WS signals, I2S slave starts transmitting data.
    \end{itemize}
    \item Wait for the interrupt signals set in Step \ref{other:i2s-item-tx-enable-corresponding-interrupts}, or check whether the transfer is completed by querying \hyperref[fielddesc:I2STXIDLE]{I2S\_TX\_IDLE}:
    \begin{itemize}
        \item 0: Transmitter is working.
        \item 1: Transmitter is in idle state.
    \end{itemize}
    \item Clear \hyperref[fielddesc:I2STXSTART]{I2S\_TX\_START} to stop data transfer.
\end{enumerate}

\subsection{Configure I2S as RX Mode}\label{sec:i2s-configure-i2s-rx-mode}

Follow the steps below to configure I2S as RX mode via software:

\begin{enumerate}
    \item Configure the clock as described in Section \ref{sec:i2s-tx-rx-clk}.
    \item Configure signal pins according to Table \ref{tab:i2s-signal-descr}.
    \item Select the mode needed by configuring \hyperref[fielddesc:I2SRXSLAVEMOD]{I2S\_RX\_SLAVE\_MOD}.
    \begin{itemize}
        \item 0: Master RX mode
        \item 1: Slave RX mode
    \end{itemize}
    \item Set needed RX data mode and RX channel mode as described in Section \ref{sec:i2s-rx-data}, and then set \hyperref[fielddesc:I2SRXUPDATE]{I2S\_RX\_UPDATE}.
    \item Reset the RX unit and its FIFO according to Section \ref{sec:i2s-reset}.
    \item Enable the corresponding interrupts. See Section \ref{sec:i2s-interrupts}.\label{other:i2s-item-rx-enable-corresponding-interrupts}
    \item Configure the GDMA receive linked list and set the receive-data length in \hyperref[regdesc:I2SRXEOFNUMREG]{I2S\_RXEOF\_NUM\_REG}.
    \item Start receiving data:
    \begin{itemize}
        \item In master mode, when the slave is ready, set \hyperref[fielddesc:I2SRXSTART]{I2S\_RX\_START} to start receiving data.
        \item In slave mode, set  \hyperref[fielddesc:I2SRXSTART]{I2S\_RX\_START} to start receiving data when BCK and WS signals are received from the master.
    \end{itemize}
    \item The received data is then stored to the specified memory address according to the configuration of GDMA. Then the corresponding interrupt set in step \ref{other:i2s-item-rx-enable-corresponding-interrupts} is generated.
\end{enumerate}


% >>> If your module has no registers, delete sample text
%       for Sections Register Summary and Registers below <<<

%%%%%%%%%%%%%%%%%%%%%%%%
%%% Register Summary %%%
%%%%%%%%%%%%%%%%%%%%%%%%

% >>> Update the sample text below and insert
%     the info on register summary into the subfile <<<

\newpage
\hypertarget{i2s-reg-summ}{}
\section{Register Summary}
\label{sec:i2s-reg-summ}

All addresses in this section are offsets relative to the I2S\regindex{n} base address. See Table \ref{tab:sysmem-base-address} \textit{\nameref{tab:sysmem-base-address}} in Chapter \ref{mod:sysmem} \textit{\nameref{mod:sysmem}} for the base addresses.

The abbreviations given in Column \textbf{Access} are explained in Section \hyperref[glossary-access-types]{\textit{Access Types for Registers}}.

\subfile{\modulefiles/45-I2S-reg-summ__EN}

%%%%%%%%%%%%%%%%%
%%% Registers %%%
%%%%%%%%%%%%%%%%%

% >>> Update the sample text below and insert
%      the info on registers into the subfile <<<

\newpage
\section{Registers}

All addresses in this section are offsets relative to the I2S\regindex{n} base address. See Table \ref{tab:sysmem-base-address} \textit{\nameref{tab:sysmem-base-address}} in Chapter \ref{mod:sysmem} \textit{\nameref{mod:sysmem}} for the base addresses.

For how to program reserved fields, please refer to Section \hyperref[programming-reserved-field]{\textit{Programming Reserved Register Field}}.

\subfile{\modulefiles/45-I2S-reg__EN}

\end{document}
